% Generated draft: requires mathematical and source review.
\documentclass[11pt]{article}
\usepackage[margin=25mm]{geometry}
\usepackage{fontspec}
\setmainfont{Latin Modern Roman}
\usepackage{amsmath,amssymb,mathtools,unicode-math}
\setmathfont{Latin Modern Math}
\usepackage{enumitem,needspace,xurl,hyperref}
\setlength{\parindent}{0pt}
\setlength{\parskip}{5pt}
\setlength{\emergencystretch}{3em}
\newcommand{\N}{\mathbb{N}}
\newcommand{\Z}{\mathbb{Z}}
\newcommand{\Q}{\mathbb{Q}}
\newcommand{\R}{\mathbb{R}}
\newcommand{\C}{\mathbb{C}}
\newcommand{\F}{\mathbb{F}}
\newcommand{\A}{\mathbb{A}}
\newcommand{\PP}{\mathbb{P}}
\newcommand{\E}{\mathbb{E}}
\newcommand{\eps}{\varepsilon}
\newcommand{\dd}{\,\mathrm{d}}
\begin{document}
\section*{UnsolvedMath ID: 701040}
\textbf{Source problem code:} catalogue entry
\subsection{Definitions and mathematical statement}
Let \(p\) be a prime and let \(\F_p=\Z/p\Z\). Under the standard interpretation of the Nottingham group, put
\[
J=N(\F_p)=\left\{f(t)=t+\sum_{n=2}^{\infty}a_nt^n : a_n\in\F_p\right\},
\]
with group operation given by composition of formal power series. Equip \(J\) with the topology whose neighbourhood basis at the identity consists of
\[
J_m=\{f\in J:f(t)\equiv t\pmod{t^{m+1}}\},\qquad m\geq1.
\]
This is the usual pro-\(p\) topology on the Nottingham group.

A pro-\(p\) group \(G\) is finitely presented as a pro-\(p\) group if there are integers \(d,r\geq0\) and elements \(\rho_1,\ldots,\rho_r\) of the free pro-\(p\) group \(F_d\) on \(d\) generators such that
\[
G\cong F_d/\overline{\langle\!\langle\rho_1,\ldots,\rho_r\rangle\!\rangle}
\]
as topological groups, where the denominator is the smallest closed normal subgroup of \(F_d\) containing all the \(\rho_i\).

The remaining assertion, part (b) of the supplied source, is that \(N(\F_2)\) is finitely presented. With finite presentation understood in the pro-\(2\) sense just defined, the requested assertion is therefore
\[
\exists\,d,r\geq0\;\exists\,\rho_1,\ldots,\rho_r\in F_d:
\quad N(\F_2)\cong F_d/\overline{\langle\!\langle\rho_1,\ldots,\rho_r\rangle\!\rangle}.
\]
Here \(F_d\) is free pro-\(2\), and the isomorphism is continuous with continuous inverse.

This is a provisional explicit interpretation: the supplied excerpt refers to Archive problem 12.24 for the definition of \(N(\Z/p\Z)\), but that definition is absent. It also does not explicitly specify whether finite presentation is meant in the pro-\(p\) or abstract group category. The definitions and category used above require confirmation against the referenced source; abstract finite presentation would be a different assertion.

For context, archived part (a) asks for finite presentation of \(N(\F_p)\) for every prime \(p>2\). The supplied archive reports this part as proved by M. V. Ershov. That report is preserved as supplied attribution, not as an independently verified theorem or literature assessment.
\subsection{Short English statement}
Kourovka problem 14.55(b), proposed by C. R. Leedham-Green, asks whether the Nottingham group over \(\F_2\) is finitely presented. The displayed formulation uses finite presentation in the pro-\(2\) category, subject to checking the omitted source definition and intended category. The supplied record labels the problem partially solved and reports the cases \(p>2\) as solved by M. V. Ershov.
\subsection{Sources}
\begin{itemize}
\item Local UnsolvedMath registry: permanent ID \(701040\), title \emph{Kourovka Notebook Problem 14.55}, exact source code \texttt{catalogue entry}. The supplied record identifies C. R. Leedham-Green as proposer and gives the proposed year as 1999.
\item \emph{The Kourovka Notebook (Unsolved Problems in Group Theory)}, 14th issue (1999), problem 14.55; supplied background identifies an update in the 21st issue (2026), arXiv:1401.0300v48. The supplied collection description names E. I. Khukhro and V. D. Mazurov as editors. Part (b) concerns \(p=2\); archived part (a) concerns \(p>2\).
\item Archive problem 12.24 is cited in the supplied statement for the definition of the Nottingham group. Its text is not supplied, so the definition and presentation category have not been checked against that reference.
\item Supplied Archive of Solved Problems attribution for \(p>2\): M. V. Ershov, \emph{J. London Math. Soc.} \textbf{71} (2005), 362--378. No article title or link is supplied.
\item Supplied status review dated \texttt{2026-10-16}: canonical status \texttt{partially\_solved}, evidence label \texttt{published}. Its notes report that the \(p>2\) part is solved in the Archive and that the remaining part is in the main section. This is supplied dated curation, not an independently verified assessment. The supplied record was also updated on \texttt{2026-10-16}.
\end{itemize}
\end{document}
